\documentclass{article}
\usepackage[utf8]{inputenc}
\usepackage[T1]{fontenc}
\usepackage{geometry}
\usepackage{graphicx}
\usepackage{float}
\usepackage{hyperref}
\usepackage[english,indonesian]{babel}

\title{Optimasi Pemrosesan Template LaTeX Menggunakan Pendekatan Berbasis AI}
\author{Academic Support Agent}
\date{\today}

\begin{document}
\maketitle

\begin{abstract}

\end{abstract}

\chapter{PENDAHULUAN}


\chapter{TINJAUAN PUSTAKA}


% Original content (commented out):
% \documentclass[12pt]{article}
\setlength{\oddsidemargin}{0in}
\setlength{\evensidemargin}{0in}
\setlength{\textwidth}{6.5in}
\setlength{\parindent}{0in}
\setlength{\parskip}{\baselineskip}

\usepackage{amsmath,amsfonts,amssymb}

\title{Your Project Title}

\begin{document}

CSCI 3104 Spring 2018\hfill Problem Set \#\\
Firstname Lastname (MM/DD)

\hrulefill

\begin{enumerate}
\item Question 1
  \begin{enumerate}
  \item Part A

  SOLUTION

  \item Part B

  SOLUTION
  \end{enumerate}
\newpage

\item Question 2
  \begin{enumerate}
  \item Part A

  SOLUTION

  \item Part B

  SOLUTION
  \end{enumerate}
\newpage

\item Question 3
  \begin{enumerate}
  \item Part A

  SOLUTION

  \item Part B

  SOLUTION
  \end{enumerate}
\newpage

\item Question 4
  \begin{enumerate}
  \item Part A

  SOLUTION

  \item Part B

  SOLUTION
  \end{enumerate}
\newpage

\item Question 5
  \begin{enumerate}
  \item Part A

  SOLUTION

  \item Part B

  SOLUTION
  \end{enumerate}
\newpage

\end{enumerate}
\end{document}

\end{document}